% \iffalse meta-comment
% 
% This is file `ltcaption.dtx'.
% 
% Copyright (C) 2007 Axel Sommerfeldt (caption@sommerfee.de)
% 
% --------------------------------------------------------------------------
% 
% This work may be distributed and/or modified under the
% conditions of the LaTeX Project Public License, either version 1.3
% of this license or (at your option) any later version.
% The latest version of this license is in
%   http://www.latex-project.org/lppl.txt
% and version 1.3 or later is part of all distributions of LaTeX
% version 2003/12/01 or later.
% 
% This work has the LPPL maintenance status "maintained".
% 
% This Current Maintainer of this work is Axel Sommerfeldt.
% 
% This work consists of the files ltcaption.ins and ltcaption.dtx,
% and the derived file ltcaption.sty.
% 
% \fi
% \CheckSum{191}
%
% \iffalse
%<*driver>
\NeedsTeXFormat{LaTeX2e}[1994/12/01]
\documentclass{ltxdoc}
\setlength\parindent{0pt}
\setlength\parskip{\smallskipamount}
%
\ifx\pdfoutput\undefined\else
  \ifcase\pdfoutput\else
    \usepackage{mathptmx,courier}
    \usepackage[scaled=0.90]{helvet}
    \addtolength\marginparwidth{15pt}
  \fi
\fi
%
\usepackage{hyperref,longtable}
\makeatletter
\let\LT@makecaption@ORI\LT@makecaption
\usepackage{ltcaption}[2007/04/15]
%
\DeclareRobustCommand{\KOMAScript}{\textsf{K\kern.05em O\kern.05em%
    M\kern.05em A\kern.1em-\kern.1em Script}}
%
\begin{document}
  \DocInput{ltcaption.dtx}
\end{document}
%</driver>
% \fi
%
% \newcommand*\purerm[1]{{\upshape\mdseries\rmfamily #1}}
% \newcommand*\puresf[1]{{\upshape\mdseries\sffamily #1}}
% \newcommand*\purett[1]{{\upshape\mdseries\ttfamily #1}}
% \let\package\puresf
% \let\env\purett \let\opt\purett
%
% \def\thispackage{the \package{ltcaption} package}
% \def\Thispackage{The \package{ltcaption} package}
% \newcommand*\version[2][]{$v#2$}
% 
% \GetFileInfo{ltcaption.sty}
% \title{\Thispackage\thanks{%
%        This package has version number \fileversion, last revised \filedate.}}
% \author{Axel Sommerfeldt\\
%         \href{mailto:caption@sommerfee.de}{\texttt{caption@sommerfee.de}}}
% \date{2007/04/15}
% \maketitle
% 
% \begin{abstract}
% This package fixes caption problems with other-than-centered aligned longtables.
% (solves \LaTeX\ PR |tools/3387|)
% \end{abstract}
% 
% \tableofcontents
%
% \clearpage
% \section{The user interface}
%
% Just include this package \emph{after} the longtable package, e.g.:
% \begin{quote}
%   |\usepackage{longtable,ltcaption}|
% \end{quote}
%
% There are two possibilities for adaption:
%
% You can alter the centering of the caption box (of width |\LTcapwidth|)
% by setting the lengths |\LTcapleft| \& |\LTcapright| to appropriate values.
% These are set to |\fill| by default, just like the values |\LTleft| \& |\LTright|.
%
% Another option is the usage of the command |\LTcapmarginsfalse| which makes
% \thispackage\ to use the values |\LTleft| \& |\LTright| instead of
% |\LTcapleft| \& |\LTcapright|.
%
% \emph{Note:}
% These adaptions do not work when used with one of the \KOMAScript\ classes
% |scrartcl|, |scrreprt| or |scrbook|, the \KOMAScript\ settings for captions
% are used instead. Same with the \package{caption} package.
%
% \clearpage
% \section{Spot the difference}
%
% \begin{minipage}{\linewidth}
% Without \thispackage:
% \makeatletter
% \let\LT@makecaption\LT@makecaption@ORI
% \makeatother
%
% \iffalse show textwidth \fi
% \noindent\rule{\textwidth}{1pt}
%
% \begin{longtable}[l]{l}
% \caption{Left aligned longtable left aligned longtable left aligned
%   longtable}\\
% This is only a test\\
% \end{longtable}
% 
% \begin{longtable}[r]{l}
% \caption{Right aligned longtable right aligned longtable right aligned
%   longtable}\\
% This is only a test\\
% \end{longtable}
% 
% \begin{longtable}[c]{l}
% \caption{Centered longtable centered longtable centered longtable
%   centered longtable}\\
% This is only a test\\
% \end{longtable}
%
% \noindent\rule{\textwidth}{1pt}
% \end{minipage}
%
% \bigskip
%
% \begin{minipage}{\linewidth}
% With \thispackage\ (and the default value of |\LTcapwidth|):
%
% \noindent\rule{\textwidth}{1pt}
%
% \begin{longtable}[l]{l}
% \caption{Left aligned longtable left aligned longtable left aligned
%   longtable}\\
% This is only a test\\
% \end{longtable}
% 
% \begin{longtable}[r]{l}
% \caption{Right aligned longtable right aligned longtable right aligned
%   longtable}\\
% This is only a test\\
% \end{longtable}
% 
% \begin{longtable}[c]{l}
% \caption{Centered longtable centered longtable centered longtable
%   centered longtable}\\
% This is only a test\\
% \end{longtable}
%
% \noindent\rule{\textwidth}{1pt}
% \end{minipage}
%
% \bigskip
%
% \begin{minipage}{\linewidth}
% With \thispackage\ and |\LTcapwidth=\linewidth|:
% \LTcapwidth=\linewidth
%
% \noindent\rule{\textwidth}{1pt}
%
% \begin{longtable}[l]{l}
% \caption{Left aligned longtable left aligned longtable
%   left aligned longtable left aligned longtable}\\
% This is only a test\\
% \end{longtable}
% 
% \begin{longtable}[r]{l}
% \caption{Right aligned longtable right aligned longtable
%   right aligned longtable right aligned longtable}\\
% This is only a test\\
% \end{longtable}
% 
% \begin{longtable}[c]{l}
% \caption{Centered longtable centered longtable centered longtable
%   centered longtable}\\
% This is only a test\\
% \end{longtable}
%
% \noindent\rule{\textwidth}{1pt}
% \end{minipage}
%
% \bigskip
%
% \begin{minipage}{\linewidth}
% With \thispackage\ and |\LTcapleft=0pt|
% resp. |\LTcapright=0pt|:
%
% \noindent\rule{\textwidth}{1pt}
%
% \LTcapleft=0pt\relax
% \LTcapright=\fill
% \begin{longtable}[l]{l}
% \caption{Left aligned longtable left aligned longtable left aligned
%   longtable}\\
% This is only a test\\
% \end{longtable}
%
% \LTcapleft=\fill
% \LTcapright=0pt\relax
% \begin{longtable}[r]{l}
% \caption{Right aligned longtable right aligned longtable right aligned
%   longtable}\\
% This is only a test\\
% \end{longtable}
%
% \noindent\rule{\textwidth}{1pt}
% \end{minipage}
%
% \bigskip
%
% \begin{minipage}{\linewidth}
% With \thispackage\ and |\LTcapleft=\tabcolsep|\\
% resp. |\LTcapright=\tabcolsep|:
%
% \noindent\rule{\textwidth}{1pt}
%
% \LTcapleft=\tabcolsep
% \LTcapright=\fill
% \begin{longtable}[l]{l}
% \caption{Left aligned longtable left aligned longtable left aligned
%   longtable}\\
% This is only a test\\
% \end{longtable}
%
% \LTcapleft=\fill
% \LTcapright=\tabcolsep
% \begin{longtable}[r]{l}
% \caption{Right aligned longtable right aligned longtable right aligned
%   longtable}\\
% This is only a test\\
% \end{longtable}
%
% \noindent\rule{\textwidth}{1pt}
% \end{minipage}
%
% \bigskip
%
% \begin{minipage}{\linewidth}
% With \thispackage\ and |\LTcapmarginsfalse|:
% \LTcapmarginsfalse
%
% \noindent\rule{\textwidth}{1pt}
%
% \begin{longtable}[l]{l}
% \caption{Left aligned longtable left aligned longtable left aligned
%   longtable}\\
% This is only a test\\
% \end{longtable}
%
% \begin{longtable}[r]{l}
% \caption{Right aligned longtable right aligned longtable right aligned
%   longtable}\\
% This is only a test\\
% \end{longtable}
%
% \begin{longtable}[c]{l}
% \caption{Centered longtable centered longtable centered longtable
%   centered longtable}\\
% This is only a test\\
% \end{longtable}
%
% \noindent\rule{\textwidth}{1pt}
% \end{minipage}
%
% \StopEventually{}
% \clearpage
% 
% \iffalse
% --------------------------------------------------------------------------- %
% \fi
% 
% \DoNotIndex{\\,\_,\ ,\@@par}
% \DoNotIndex{\@classoptionslist,\@currext,\@currname}
% \DoNotIndex{\@ehc,\@ehd,\@empty,\@expandtwoargs}
% \DoNotIndex{\@for,\@firstofone,\@firstoftwo}
% \DoNotIndex{\@gobble,\@gobblefour,\@gobbletwo,\@hangfrom}
% \DoNotIndex{\@ifnextchar,\@ifstar,\@ifundefined,\@latex@error}
% \DoNotIndex{\@namedef,\@nameuse}
% \DoNotIndex{\@onlypreamble,\@parboxrestore,\@plus,\@ptionlist}
% \DoNotIndex{\@removeelement,\@restorepar,\@secondoftwo,\@setpar}
% \DoNotIndex{\@tempa,\@tempboxa,\@tempdima,\@tempdimb,\@tempdimc,\@tempb,\@tempc}
% \DoNotIndex{\@testopt}
% \DoNotIndex{\@undefined,\@unprocessedoptions,\@unusedoptionlist}
% \DoNotIndex{\p@,\z@}
% \DoNotIndex{\active,\addtocounter,\addtolength,\advance}
% \DoNotIndex{\baselineskip,\begin,\begingroup,\bfseries,\box}
% \DoNotIndex{\catcode,\centering,\changes,\csname,\def,\divide,\do,\downarrow}
% \DoNotIndex{\edef,\else,\empty,\end,\endcsname,\endgraf,\endgroup,\expandafter}
% \DoNotIndex{\fi,\footnotesize,\global}
% \DoNotIndex{\hangindent,\hbox,\hfil,\hsize,\hskip,\hspace,\hss}
% \DoNotIndex{\ifcase,\ifdim,\ifnum,\ifodd,\ifvoid,\ifvmode}
% \DoNotIndex{\ifx,\ignorespaces,\itshape}
% \DoNotIndex{\Large,\large,\leavevmode,\leftmargini,\leftskip,\let,\linewidth}
% \DoNotIndex{\llap,\long,\m@ne,\margin,\mdseries,\message}
% \DoNotIndex{\newcommand,\newdimen,\newlength,\newline,\newif,\newsavebox}
% \DoNotIndex{\next,\nobreak,\nobreakspace,\noexpand,\noindent,\numberline}
% \DoNotIndex{\normalsize,\or,\par,\parbox,\parfillskip}
% \DoNotIndex{\parindent,\parskip,\prevdepth,\protect,\protected@edef,\protected@write}
% \DoNotIndex{\providecommand,\quad}
% \DoNotIndex{\raggedleft,\raggedright,\relax,\renewcommand,\RequirePackage}
% \DoNotIndex{\rightskip,\rmfamily}
% \DoNotIndex{\sbox,\scriptsize,\scshape,\setbox,\setlength,\sffamily,\slshape}
% \DoNotIndex{\small,\string,\space,\strut}
% \DoNotIndex{\textheight,\the,\toks@,\typeout,\ttfamily}
% \DoNotIndex{\unvbox,\uparrow,\upshape,\usebox,\usepackage}
% \DoNotIndex{\value,\vbox,\vsize,\vskip,\wd,\width,\z@skip}
% \DoNotIndex{\AtBeginDocument,\AtEndOfPackage,\CurrentOption,\DeclareOption}
% \DoNotIndex{\ExecuteOptions,\GenericWarning,\IfFileExists,\InputIfFileExists}
% \DoNotIndex{\NeedsTeXFormat,\MessageBreak}
% \DoNotIndex{\PackageError,\PackageInfo,\PackageWarning,\PackageWarningNoLine}
% \DoNotIndex{\ProcessOptions,\ProvidesPackage}
%
% \iffalse
% --------------------------------------------------------------------------- %
% \fi
% 
% \setlength{\parskip}{0pt plus 1pt}
%
% \changes{v1.0}{2007/04/15}{First release}
%
% \iffalse
% --------------------------------------------------------------------------- %
% \fi
%
% \clearpage
% \section{The Implementation}
% \iffalse
%<*package>
% \fi
%
% \subsection{Identification}
%
%    \begin{macrocode}
\NeedsTeXFormat{LaTeX2e}[1994/12/01]
\ProvidesPackage{ltcaption}[2007/04/15 v1.0 longtable captions (AR)]
%    \end{macrocode}
%
% \subsection{User interface}
%
% \begin{macro}{\LTcapleft}
% \begin{macro}{\LTcapright}
% \begin{macro}{\ifLTcapmargins}
% Our skips and the flag belonging to them.
%    \begin{macrocode}
\newskip\LTcapleft \LTcapleft=\fill
\newskip\LTcapright \LTcapright=\fill
\newif\ifLTcapmargins \LTcapmarginstrue
  % default: use these skips (and not \LTleft & \LTright)
%    \end{macrocode}
% \end{macro}
% \end{macro}
% \end{macro}
%
% \subsection{The longtable patch}
%
% \begin{macro}{\LT@makecaption}
%  |\LT@makecaption|\marg{cmd}\marg{label}\marg{text}\par
% \smallskip
%  Original code:
%  \begin{verbatim}
%  \def\LT@makecaption#1#2#3{%
%    \LT@mcol\LT@cols c{\hbox to\z@{\hss\parbox[t]\LTcapwidth{%
%      % Based on article class "\@makecaption", "#1" is "\@gobble" in star
%      % form, and "\@firstofone" otherwise.
%      \sbox\@tempboxa{#1{#2: }#3}%
%      \ifdim\wd\@tempboxa>\hsize
%        #1{#2: }#3%
%      \else
%        \hbox to\hsize{\hfil\box\@tempboxa\hfil}%
%      \fi
%      \endgraf\vskip\baselineskip}%
%    \hss}}}
%  \end{verbatim}
%  Our code:
%    \begin{macrocode}
\renewcommand*\LT@makecaption[3]{%
  \caption@LT@make{\hb@xt@\hsize{%
    \ifLTcapmargins
      \hspace\LTcapleft
    \else
      \hspace\LTleft
    \fi
    \parbox[t]\LTcapwidth{%
      \sbox\@tempboxa{#1{#2: }#3}%
      \ifdim\wd\@tempboxa>\hsize
        #1{#2: }#3%
      \else
        \hbox to\hsize{\hfil\box\@tempboxa\hfil}%
      \fi
      \endgraf\vskip\baselineskip}%
    \ifLTcapmargins
      \hskip\LTcapright
    \else
      \hskip\LTright
    \fi}}}
%    \end{macrocode}
% \end{macro}
%
% \begin{macro}{\caption@LT@make}
%  Typesets the caption as |\multicolumn|\ldots
%    \begin{macrocode}
\newcommand\caption@LT@make[1]{%
  \noalign{\caption@LT@config}%
%    \end{macrocode}
% Note: If used with the \package{array} package |\caption@LTfmt| needs
% to be expanded, therefore we need some |\expandafter| here.
%    \begin{macrocode}
  \expandafter\LT@mcol\expandafter\LT@cols\expandafter{\caption@LTfmt}{%
    \hb@xt@\z@{%
      \hspace\caption@LTleft
      \parbox[t]\linewidth{#1}%
      \hspace\caption@LTright}}}%
%    \end{macrocode}
% \end{macro}
%
% \begin{macro}{\caption@LT@config}
% |\caption@LT@config| analyses |\LTleft| \& |\LTright|
% and set |\caption@LTleft| \& |\caption@LTright| accordingly
% to the `opposite' values, e.g., |\LTleft=1cm| will result to
% |\caption@LTleft=-1cm| and |\LTleft=0pt plus 1fill|
% will result to |\caption@LTleft=0pt minus 1fill|.
% Furthermore |\caption@LTfmt| is set to the according
% multicolumn format; this is far away from being bulletproof
% (e.g., a stretch or shrink will always be treated as `fill')
% but will hopefully cover all `real' cases.
%    \begin{macrocode}
\newcommand*\caption@LT@config{%
%    \end{macrocode}
%    \begin{macrocode}
  \caption@LT@parse\LTleft\caption@LTleft\caption@ifLTleft
  \caption@LT@parse\LTright\caption@LTright\caption@ifLTright
%    \end{macrocode}
%    \begin{macrocode}
  \xdef\caption@LTfmt{%
    @{}\caption@ifLTleft{\caption@ifLTright{c}{r}}{l}@{}}}
%    \end{macrocode}
% \end{macro}
%
%    \iffalse
% \def\@defplusminus#1 #2 #3 #4 #5\@nil{%
%   \def\@plus{#2}\def\@minus{#4}}
% \@tempskipa=1pt plus 1pt minus 1pt\relax
% \expandafter\@defplusminus\the\@tempskipa\@nil
%    \fi
%
% \begin{macro}{\caption@LT@parse}
% Parsing of the skip, we collect a |\@fixpart|, a |@pluspart|,
% and a |\@minuspart| and make our definitions based on that.
%    \begin{macrocode}
\newcommand*\caption@LT@parse[3]{%
%    \end{macrocode}
%    \begin{macrocode}
  \let\@pluspart\@undefined
  \let\@minuspart\@undefined
  \expandafter\caption@LT@parse@\expandafter\@fixpart\the#1 x %
%    \end{macrocode}
%    \begin{macrocode}
  \xdef#2{-\@fixpart
    \ifx\@minuspart\@undefined\else
      \space\@plus\space\@minuspart
    \fi
    \ifx\@pluspart\@undefined\else
      \space\@minus\space\@pluspart
    \fi}%
%    \end{macrocode}
%    \begin{macrocode}
  \ifx\@pluspart\@undefined
    \ifx\@minuspart\@undefined
      \let#3\@secondoftwo
    \else
      \let#3\@firstoftwo
    \fi
  \else
    \let#3\@firstoftwo
  \fi}
%    \end{macrocode}
%  Note: |\def\@tempa{#2}\ifx\@tempa\@plus| does not work
%  because of different catcodes.
%    \begin{macrocode}
\def\caption@LT@parse@#1#2 {%
  \edef\@tempa{\@car#2\@nil}%
  \if p\@tempa
    \def\next{\caption@LT@parse@\@pluspart}%
  \else\if m\@tempa
    \def\next{\caption@LT@parse@\@minuspart}%
  \else\if x\@tempa
    \let\next\relax
  \else
    \def#1{#2}%
    \def\next{\caption@LT@parse@ @}%
  \fi\fi\fi
  \next}
%    \end{macrocode}
% \end{macro}
%
% \subsection{Adaption for \KOMAScript}
%
%    \begin{macrocode}
\@ifundefined{@komalongtablefalse}{}{%
  \if@komalongtable
    \renewcommand{\LT@makecaption}[3]{%
      \noalign{%
        \if@captionabove
          \vskip\belowcaptionskip
        \else
          \vskip\abovecaptionskip
        \fi
      }%
      \caption@LT@make{%
        \@@makecaption{#1}{#2}{#3}%
        \endgraf
        \if@captionabove
          \vskip\abovecaptionskip
        \else
          \vskip\belowcaptionskip
        \fi
      }%
    }%
  \fi}
%    \end{macrocode}
%
% \iffalse
%</package>
% \fi
%
% \iffalse
% --------------------------------------------------------------------------- %
% \fi
%
% \begin{thebibliography}{9}
%   \iffalse
%     TODO: \href{...}{...}
%   \fi
%   \bibitem{longtable}
%   David Carlisle:
%   \emph{The longtable package},
%   2004/02/01
% \end{thebibliography}
%
% \iffalse
% --------------------------------------------------------------------------- %
% \fi
%
% \clearpage
% \Finale
%
\endinput
